\section{Threat Model}
\label{sec:threat}

\paragraph{Setting.}
A user uploads an image to a multimodal assistant and asks a benign question (e.g.\ ``describe this image'' or ``extract all text from this screenshot''). The attacker controls only the image pixels; they cannot edit the user's prompt, the system prompt, or the model weights. The attacker chooses one short \emph{target phrase} (a URL, a payment-information request, a piece of misinformation, etc.) before the attack begins, and succeeds if that phrase --- or its semantic class --- appears in the assistant's reply.

\paragraph{Capabilities.}
We grant the attacker:
\begin{itemize}
  \item \textbf{Perfect knowledge of a small white-box ensemble.} The attacker can compute gradients through two to four open \vlm{}s (Qwen2.5-VL-3B, Qwen2-VL-2B, DeepSeek-VL-1.3B, BLIP-2-OPT-2.7B). These are the ``surrogates''.
  \item \textbf{One target phrase per attack.} The optimisation objective is fixed at attack-construction time and reused across many test images.
  \item \textbf{No access to the deployed (black-box) target.} The attacker does not query the deployed model during optimisation; transferability is what we measure in \S\ref{sec:discussion}.
\end{itemize}

\paragraph{Constraints.}
Two constraints make the attack realistic rather than merely possible:
\begin{itemize}
  \item \textbf{Perceptual budget.} The adversarial image must remain visually indistinguishable from the clean photo. We enforce $\Linf \leq 16/255$, which yields PSNR $\approx 25.2$~dB on every test image we use \citep{goodfellow2015explaining}.
  \item \textbf{Benign user prompt.} The user does not collaborate with the attacker. The reply is judged against a clean baseline answer to the same prompt on the same image (with the noise removed).
\end{itemize}

\paragraph{Attack scenarios.}
This setting is not hypothetical. We highlight three operational scenarios:
\begin{enumerate}
  \item \textbf{Hosted assistant uploads.} A user shares a screenshot or photo with a chat assistant. The model's reply is rendered to the user --- so any embedded URL, phone number, or payment instruction in the reply becomes a clickable target inside a trusted UI.
  \item \textbf{LLM agents reading screenshots.} Browser-controlling agents observe the live page through screenshots. A poisoned page or banner ad can therefore steer what the agent perceives and what it does next.
  \item \textbf{Tool replays / MCP.} Tool outputs (image attachments, OCR responses, retrieved documents) are fed back into the model. A persistent adversarial image inside a tool response can re-inject the payload across many turns of a long-running agent.
\end{enumerate}

The common thread is that the model treats user-supplied imagery as ground truth: it is the cheapest form of input to forge, and every defence below the model itself --- mailbox filters, link checkers, prompt sanitisers --- runs after the visual channel has already been read.
